\scanpage{27}
En effet, de la suite exacte
\[
0\unsure{-|\to} A\xrightarrow{x}A\to A/xA\to0
\]
on tire
\[
H^{n-1}(A)\xrightarrow{x}H^{n-1}(A)\to H^{n-1}(A/xA)\to H^n(A)\xrightarrow{x}H^n(A)
\]
et compte tenu de l'hypothèse
\[
H^{n-1}(A/xA)\simeq\deleted{\operatorname{Ker}\text{ de }x\text{ dans}}\inserted{\text{ann. de }x}\text{ dans }H^n(A)
\]
\deleted{\unsure{a}}
Comme $H^n(A)\simeq D(\Omega)$, on trouve bien
\[
H^{n-1}(A/xA)\simeq D(\Omega/x\Omega),\qquad\text{cqfd.}
\]

\uline{Corollaire} Supposons $A$ Coh. Mac. et soit
$x_1\mathord{-}\,\mathord{-}x_p$ une partie d'un syst. de
paramètres de $A$. Alors $\Omega_A/(x_1,\mathord{.}\,\mathord{-},x_p)\Omega_A$ est
un module fondamental pour $A/(x_1,\mathord{-},x_p)$.

Supposons maintenant que $A$ soit un
quotient d'un anneau local \deleted{\uline{\unsure{régulier}}}\inserted{Coh. Mac.} $A'$
\deleted{\unsure{par}}\inserted{de dim \unsure{$n'$|$N'$}} 
\[
A\simeq A'/\unsure{\mathfrak N|N}
\]
Supposons $A'$ muni d'un $\Omega_{A'}$. Si
$M$ est un $A$-module, c'est un $A'$-module,
et $D(M)$ peut-être interprété ainsi
[choisissons un module dualisant $I'$ \inserted{$=H^{\unsure{n'|N'}}(\Omega_{A'})$} pour $A'$,
en \unsure{$=$|$\simeq$} temps que $\Omega_{A'}$]. Car tout \unsure{hom}
% archive sheet number: 26

